\documentclass[11pt]{article}

\usepackage[margin=0.85in]{geometry}
\usepackage[T1]{fontenc}
\usepackage{lmodern}
\usepackage{microtype}
\usepackage{amsmath,amssymb,amsthm,mathtools}
\usepackage{graphicx}
\usepackage{xcolor}
\usepackage[hidelinks]{hyperref}

\definecolor{statusblue}{HTML}{1F4E79}
\newtheorem{theorem}{Theorem}
\newtheorem{lemma}{Lemma}
\newcommand{\cnull}{c_0(X)}

\title{A reflexive smooth support space preserving weak convergence}
\author{Solution packet for arXiv:1505.06235}
\date{22 July 2026}

\begin{document}
\maketitle

\begin{center}
\fcolorbox{statusblue}{statusblue!5}{%
\parbox{0.92\textwidth}{
\textbf{Status: claimed full solution; likely valid.}
Every weakly convergent sequence of Borel probability measures on a separable
Banach space can be realized on one separable reflexive Banach subspace that
is compactly embedded in the original space, has an equivalent
Fr\'echet-differentiable norm, carries all the measures, and preserves their
weak convergence.}}
\end{center}

\section{Source question}

Ostrovsky and Sirota prove that weak convergence of probability measures on a
separable Banach space $X$ can be preserved after passing to a separable
Banach subspace compactly embedded in $X$ \cite{OS}.  In concluding
remark~C on page~8, they recall Buldygin's theorem that a probability measure
may be concentrated on such a subspace which is additionally reflexive and
has a Fr\'echet-differentiable norm, and ask whether one can retain these
``good'' properties while also preserving the weak convergence.  The source
passage is reproduced below.

\begin{center}
\includegraphics[width=\textwidth]{figures/open_problem_crop.png}
\end{center}

In precise terms: if $\mu_n\Rightarrow\mu$ on $X$, can one choose a single
separable reflexive Banach space $Y\subset X$, compactly embedded in $X$ and
Fr\'echet-smoothly renormable, such that every $\mu_n$ and $\mu$ is
concentrated on $Y$ and $\mu_n\Rightarrow\mu$ also as measures on $Y$?  The
answer is yes.

\section{Input theorem and proof intuition}

We use the following published support theorem of Buldygin, in exactly the
form recalled in the source paper \cite{Buldygin}: every Borel probability
measure on a separable Banach space $B$ is concentrated on the range of a
compact continuous injection $j:R\to B$, where $R$ is a separable reflexive
Banach space.  The smooth norm will be recovered at the end by equivalent
renorming.

The main difficulty is that putting all measures on one reflexive support
does not itself make the family tight in the stronger support norm.  We avoid
this by using a Skorokhod coupling $\xi_n\to\xi$ almost surely and treating the
whole difference sequence
\[
   D=(\xi_n-\xi)_{n\geq1}
\]
as one random variable with values in $\cnull$.  Buldygin's theorem puts $D$
on a reflexive space $Z$ compactly embedded in $\cnull$.  Compactness makes
the coordinate operators $Z\to X$ uniformly small in the tail.

Those coordinate ranges still need to be assembled inside one reflexive
subspace of $X$.  We group them into finite blocks, take an $\ell_1$-sum
inside each block, and an $\ell_2$-sum across blocks.  Increasing block
weights make the norm of the $n$th coordinate tend to zero, while sufficiently
small coordinate-operator tails make the assembled operator compact.  This
is the only new mechanism in the proof.

\section{Compact coordinate assembly}

For $n\geq1$, let $Q_n:\cnull\to X$ be the $n$th coordinate map.

\begin{lemma}[reflexive coordinate assembly]\label{lem:assembly}
Let $Z$ be a separable reflexive Banach space and
$J:Z\to\cnull$ a compact operator.  Set $A_n=Q_nJ$.  Then there are a
separable reflexive Banach space $\mathcal E$ and a compact operator
$T:\mathcal E\to X$ such that
\[
  A_nz\in T(\mathcal E)
  \quad\text{and}\quad
  \|A_nz\|_{T(\mathcal E)}\longrightarrow0
  \qquad(z\in Z),
\]
where the range is equipped with the quotient norm induced by $T$.
\end{lemma}

\begin{proof}
Compact subsets of $\cnull$ have uniformly vanishing coordinate tails.
Indeed, this follows by taking a finite uniform-norm net and then taking a
common tail for its finitely many centers.  Applied to the closure of
$J(B_Z)$, this gives
\begin{equation}\label{eq:tail}
  \sup_{n\geq N}\|A_n\|\longrightarrow0.
\end{equation}

Choose integers
\[
  1=N_1<N_2<N_3<\cdots
\]
so that, for $k\geq2$,
\begin{equation}\label{eq:blocksmall}
  \sup_{n\geq N_k}\|A_n\|\leq4^{-k}.
\end{equation}
Put $I_k=\{N_k,\ldots,N_{k+1}-1\}$ and
\[
  E_k=\ell_1(I_k;Z),
  \qquad
  \mathcal E=\left(\bigoplus_{k\geq1}E_k\right)_{\ell_2}.
\]
Each $I_k$ is finite, hence each $E_k$ is reflexive; consequently
$\mathcal E$ is reflexive and separable.  Define
\begin{equation}\label{eq:T}
  T((z_n)_{n\geq1})
  =\sum_{k\geq1}2^k\sum_{n\in I_k}A_nz_n.
\end{equation}
The norm of the $k$th block map is at most
\[
  \beta_k:=2^k\sup_{n\in I_k}\|A_n\|.
\]
Here $\beta_1<\infty$ and, by \eqref{eq:blocksmall},
$\beta_k\leq2^{-k}$ for $k\geq2$.  Thus $(\beta_k)\in\ell_2$, and
Cauchy--Schwarz shows that the series in \eqref{eq:T} converges and defines a
bounded operator.

Each $A_n=Q_nJ$ is compact.  Therefore every finite block truncation of $T$
is compact, while the operator norm of the omitted tail is at most
\[
  \left(\sum_{k>K}\beta_k^2\right)^{1/2}\longrightarrow0.
\]
Hence $T$ is compact.

If $n\in I_k$ and $z\in Z$, place $2^{-k}z$ in the $n$th coordinate of
$E_k$ and zero elsewhere.  Formula \eqref{eq:T} maps this vector to $A_nz$,
and its $\mathcal E$-norm is $2^{-k}\|z\|_Z$.  Consequently
\begin{equation}\label{eq:rangenorm}
  \|A_nz\|_{T(\mathcal E)}\leq2^{-k}\|z\|_Z
  \quad(n\in I_k),
\end{equation}
which tends to zero with $n$.
\end{proof}

\section{Full affirmative answer}

\begin{theorem}\label{thm:main}
Let $X$ be a separable Banach space and let $\mu_n,\mu$ be Borel probability
measures on $X$ such that $\mu_n\Rightarrow\mu$ weakly.  There exists a
separable reflexive Banach space $Y\subset X$ with the following properties:
\begin{enumerate}
\item the inclusion $Y\hookrightarrow X$ is compact;
\item $\mu_n(X\setminus Y)=\mu(X\setminus Y)=0$ for every $n$;
\item after the measures are regarded as measures on $Y$,
      $\mu_n\Rightarrow\mu$ weakly on $Y$;
\item $Y$ admits an equivalent norm which is Fr\'echet differentiable on its
      unit sphere.
\end{enumerate}
Thus the open question in concluding remark~C of \cite{OS} has an affirmative
answer.
\end{theorem}

\begin{proof}
Since $X$ is Polish, the Skorokhod representation theorem gives $X$-valued
random variables $\xi_n,\xi$ on one probability space, with laws
$\mu_n,\mu$, such that
\begin{equation}\label{eq:skorokhod}
  \xi_n\longrightarrow\xi\quad\text{almost surely in }X.
\end{equation}
After changing the variables on a null set, define the $\cnull$-valued random
variable
\[
  D=(d_n)_{n\geq1},\qquad d_n=\xi_n-\xi.
\]
It is Borel measurable because its finite coordinate truncations are
measurable and converge to it in the $c_0(X)$ norm.

Apply Buldygin's support theorem twice.  For the law of $\xi$ there are a
separable reflexive space $R$ and a compact injection $j_0:R\to X$ such that
\begin{equation}\label{eq:limitlift}
  \xi\in j_0(R)\quad\text{almost surely}.
\end{equation}
For the law of $D$ there are a separable reflexive space $Z$ and a compact
injection $J:Z\to\cnull$ such that
\begin{equation}\label{eq:difflift}
  D\in J(Z)\quad\text{almost surely}.
\end{equation}

Apply Lemma~\ref{lem:assembly} to $J$, obtaining $\mathcal E$ and $T$.
Let
\[
  \mathcal H=R\oplus_2\mathcal E,
  \qquad
  S:\mathcal H\to X,\quad S(r,e)=j_0r+Te.
\]
The space $\mathcal H$ is separable and reflexive, and $S$ is compact.  Set
\[
  Y=\mathcal H/\ker S
\]
and identify $Y$ with $S(\mathcal H)\subset X$, equipped with the quotient
norm.  Then $Y$ is separable and reflexive, and its inclusion into $X$ is
compact.

On the probability-one event where \eqref{eq:limitlift} and
\eqref{eq:difflift} hold, write $\xi=j_0r$ and $D=Jz$.  Since
$d_n=Q_nJz=A_nz$, Lemma~\ref{lem:assembly} shows that $d_n\in Y$ and
\[
  \|d_n\|_Y\longrightarrow0.
\]
Thus $\xi,\xi_n=\xi+d_n$ all belong to $Y$ almost surely and
\begin{equation}\label{eq:Yas}
  \xi_n\longrightarrow\xi\quad\text{almost surely in }Y.
\end{equation}

The continuous injection of the separable Banach space $Y$ into the Polish
space $X$ has Borel range and Borel inverse on its range by the
Lusin--Souslin theorem.  Hence the variables above are genuinely measurable
as $Y$-valued random variables, and their laws are precisely the unique lifts
of $\mu_n,\mu$ to $Y$.  For every bounded continuous $f:Y\to\mathbb R$,
\eqref{eq:Yas} and dominated convergence give
\[
  \int_Y f\,d\mu_n\longrightarrow\int_Y f\,d\mu.
\]
This proves weak convergence on $Y$.

Finally, a separable reflexive Banach space has separable dual, and every
Banach space with separable dual admits an equivalent norm which is
Fr\'echet differentiable away from the origin \cite{DGZ}.  Replace the
quotient norm of $Y$ by such an equivalent norm.  Equivalent renorming
preserves compactness of the inclusion, concentration, and convergence, and
gives property~4.
\end{proof}

\section{Verification, scope, and novelty}

The proof separates the external inputs from the new argument.  Its external
inputs are the Skorokhod representation theorem, Buldygin's single-measure
reflexive-support theorem, Lusin--Souslin measurability, and the classical
smooth-renorming theorem.  The coordinate-assembly lemma is proved in full.
The critical estimates are \eqref{eq:blocksmall}--\eqref{eq:rangenorm}; they
simultaneously give compactness of $T$ and convergence of every coordinate in
the new range norm.  No computational step is used.

A bounded novelty search on 22 July 2026 covered arXiv and web searches for
arXiv:1505.06235, the exact paper title and open-question wording, Buldygin's
support theorem and DOI, and combinations of ``compactly embedded,''
``reflexive subspace,'' ``Fr\'echet differentiable norm,'' and ``weak
convergence of probability/Radon measures.''  It found the source paper,
Buldygin's support result, and related later work with additional moment or
uniform-integrability assumptions, but no later result explicitly resolving
this exact question.  Novelty confidence is \emph{moderate}; exhaustive
priority is not claimed.

The theorem answers the source's probabilistic question for sequences.  It
does not claim the analogous statement for the source's arbitrary nets of
sigma-finite measures.  Recommended human review: \textbf{verify and
circulate}, focusing on the interpretation of Buldygin's theorem, the compact
block assembly, and the Borel lifting step.

\small
\begin{thebibliography}{9}
\setlength{\itemsep}{2pt}
\bibitem{OS}
E.~Ostrovsky and L.~Sirota,
\emph{Strengthening of weak convergence for Radon measures in separable
Banach spaces}, 2015,
\href{https://arxiv.org/abs/1505.06235}{arXiv:1505.06235}.

\bibitem{Buldygin}
V.~V.~Buldygin,
\emph{The supports of probability measures in separable Banach spaces},
Theory Probab. Appl. \textbf{29} (1985), 546--549;
\href{https://doi.org/10.1137/1129067}{doi:10.1137/1129067}.

\bibitem{DGZ}
R.~Deville, G.~Godefroy, and V.~Zizler,
\emph{Smoothness and Renormings in Banach Spaces},
Pitman Monographs and Surveys in Pure and Applied Mathematics 64,
Longman Scientific \& Technical, 1993.
\end{thebibliography}

\end{document}
